% A downward parabola with a local maximum at (p, f(p)); the tangent
% line there is horizontal, illustrating f'(p) = 0 at an interior
% extremum.
% Source: book/tex/calculus/differentiate.tex (§ Local extrema).
\documentclass[border=2mm]{standalone}
\usepackage{tikz}
\begin{document}
\begin{tikzpicture}[scale=1.1,
    dot/.style={circle,fill,inner sep=1.6pt}]
  \draw[->] (-1.3,0)--(5.4,0) node[right] {$x$};
  \draw[->] (0,-1.3)--(0,2.7) node[above] {$y$};
  \draw[blue,thick,domain=-1:5,samples=100]
    plot (\x, {2-(\x-2)*(\x-2)/3});
  % maximum at (2,2)
  \node[dot,red] at (2,2) {};
  \node[red,above] at (2,2) {$(p,f(p))$};
  % horizontal tangent
  \draw[red] (0.3,2)--(3.7,2);
\end{tikzpicture}
\end{document}
